\chapter{Results}\label{ch:results}
% WORD BUDGET: ~3000 words (12k total target; COMP0234 limit UNVERIFIED — confirm on Moodle)
% Every number in this chapter is read from a committed results/*.json artifact.

I present results in the order the evidence was built: the baseline
landscape that reframed the project (\S\ref{sec:res-maxpressure}), the
base-model factorial (\S\ref{sec:res-factorial}), the offline distillation
tables (\S\ref{sec:res-offline}), the closed-loop outcome and its
robustness checks (\S\ref{sec:res-closedloop}), the throughput guards that
determine which of those numbers are comparable (\S\ref{sec:res-guards}),
the authorship finding this dissertation headlines
(\S\ref{sec:res-authorship}), the failure forensics the audit layer makes
possible (\S\ref{sec:res-forensics}), the audit, trust and
authorisation-re-test results (\S\ref{sec:res-audit}, \S\ref{sec:res-trust},
\S\ref{sec:res-lee}), and the serving-reliability evidence
(\S\ref{sec:res-reliability}).

\section{Baseline landscape: what the fixed-time floor revealed}\label{sec:res-maxpressure}
Adding a naive fixed-time floor to every comparison was the single most
consequential rigour decision of the project, and it produced a result I
have had to correct twice.

In the exploratory sweep (three seeds per cell, daily-resolution demand)
MaxPressure --- throughput-optimal in theory \cite{varaiya2013maxpressure}
and dominant in the simulation literature --- appeared \emph{worse than
naive fixed-time on every real London topology tested} (Euston $+15\%$,
Old Street $+5\%$, Bloomsbury $+3\%$) while dominating the synthetic grids
($-27\%$ on grid4$\times$4, $-19\%$ on grid3$\times$3). That contrast was
striking enough to become a headline claim of an earlier draft.

It does not survive replication. Table~\ref{tab:mpfloor} recomputes the
comparison paired per seed over the thirty disjoint evaluation seeds, on
the same measured-demand topologies used everywhere else in this
dissertation.

\begin{table}[h]
\caption{MaxPressure vs fixed-time, paired per seed ($n{=}30$).
Positive = MaxPressure worse. CI is 95\% on the mean relative change.}
\label{tab:mpfloor}
\centering
\small
\begin{tabular}{lrrrr}
\toprule
Topology & MaxPressure & fixed-time & rel. & 95\% CI \quad ($p$) \\
\midrule
Euston peak-hour & 259.6\,s & 235.6\,s & $+19.8\%$ & $[+3.3, +36.2]$ \quad (.21) \\
Bloomsbury grid & 527.9\,s & 527.3\,s & $+0.2\%$ & $[-1.5, +2.0]$ \quad (.90) \\
Old Street & 148.5\,s & 159.4\,s & $-6.1\%$ & $[-10.8, -1.3]$ \quad (\textbf{.016}) \\
\bottomrule
\end{tabular}
\end{table}

At $n{=}30$ the picture is not a reversal but a \finding{topology-dependent
draw}. On Old Street --- the most complex geometry in the study, a five-arm
junction with seventeen movements --- MaxPressure is \emph{significantly
better} than fixed-time, contradicting the exploratory result outright. On
Bloomsbury the two are indistinguishable. Only on the saturated Euston
corridor does MaxPressure remain worse in the mean, and there the paired
test does not reach significance ($p{=}.21$) because the per-seed spread is
large; note also that the confidence interval on the mean relative change
excludes zero while the paired test on absolute delay does not --- a
divergence caused by skew, and a reminder that on this topology a mean is a
fragile summary.

I therefore retract the general claim. What the evidence supports is
narrower and still worth stating: MaxPressure's advantage over a naive
fixed-time plan is \emph{not} reliable on irregular real geometry under
measured demand --- it is real on one topology, absent on another, and
negative-but-uncertain on a third --- whereas on the regular synthetic
grids the literature favours it is uniformly and substantially better.
The backpressure theory offers a mechanism for why the arterial is its
weakest case: the throughput-optimality proof assumes unbounded queue
storage, and finite-capacity treatments prove work-conservation losses
\cite{gregoire2014backpressure}, while the only stability results for a
deployable variant hold under heavy load \cite{zaidi2014backpressure} ---
conditions a saturated short-link corridor violates in both directions. I
offer that as a hypothesis consistent with the Euston direction, not as a
demonstrated explanation, since the effect it would explain is itself not
significant.

The practical consequence for this dissertation is unchanged and is the
reason the floor was added: any SLM-vs-MaxPressure comparison on real
topologies is confounded by MaxPressure's own variable strength, so every
controller claim here is stated against the fixed-time floor, with
MaxPressure reported alongside.

\section{Base-model factorial}\label{sec:res-factorial}
The base-model sweep (eight servable catalogue models --- the 6\,GB card
caps the ladder at $\sim$4B, with 7B-class OOM failures and a newer
generation blocked by a runtime schema change both recorded --- $\times$
four prompt formats $\times$ six topologies) established three facts
before any fine-tuning. The design is deliberately unbalanced: three
models were swept only on the original topologies and never on the later
measured-demand corridor, and phi-4-mini-reasoning is latency-excluded
(\S\ref{sec:met-baselines}) after a single probe cell. I therefore
describe it as a sweep rather than a full factorial, and the
accompanying ANOVA is fitted on the balanced subset.

First, in the exploratory sweep (three seeds per cell, daily-resolution
demand) shielded stock SLM control beat the fixed-time floor on every
topology tested --- Euston corridor $-16\%$, Bloomsbury $-2\%$, Old Street
$-14\%$ (seed-noisy), grid3$\times$3 $-21\%$, grid4$\times$4 $-28\%$. I
flag immediately that this result does \emph{not} survive replication on
the measured-demand corridor, and the powered numbers below supersede it;
the exploratory sweep is reported because it is what motivated the powered
design, not as evidence for the claim. Second, results must be conditioned
on topology: pooling across
maps is a Simpson's-paradox trap, since corridor wins and early grid
blow-ups average into a meaningless mean. Third, per-decision reasoning
quality does not predict closed-loop value: on 80 identical junction states,
cloud frontier models agree with the delay-optimal heuristic at 92--96\%
against 51--72\% for local SLMs, yet a near-chance local model still won its
closed-loop cells. The controller's value is systemic --- the SLM, shield,
and event gating together --- not raw single-junction reasoning. This
decision battery is open-loop by construction (the frontier teacher cannot be
called from inside the simulation loop) and is stated as such.

The explanatory two-way ANOVA over the factorial confirms the structure:
model and configuration interact strongly (interaction
$F{=}32$, $p{<}0.001$, $\eta^2{=}0.86$), with the effect driven by outlier
cells rather than a smooth gradient --- a further argument against pooled
claims.

The powered stock-model replication ($n{=}30$ per cell, four catalogue
models $\times$ two formats, both baselines) sharpens the picture per
topology. On \textbf{Old Street}, stock SLM control achieves the
campaign's strongest powered wins against the fixed-time floor ---
Qwen2.5-0.5B sota $-11.8\%$ ($p{<}0.001$), Qwen3-0.6B sota $-9.6\%$
($p{<}0.001$), Phi-4-mini myopic $-8.6\%$ ($p{<}0.001$) --- with one arm
also significantly beating MaxPressure (Qwen2.5-0.5B sota, $-4.9\%$,
$p{=}.014$).

On \textbf{Euston peak-hour} under the measured DfT demand profile the
result reverses, and I report it as measured: every stock arm is
\risk{7--20\% worse than the fixed-time floor} ($+7.4\%$ to $+19.7\%$),
significantly so on four of the eight arms (two-sided $p$ between .017
and .053), while remaining statistically indistinguishable from
MaxPressure. This is not a null; it is a measured harm, and it refutes
the exploratory sweep's corridor win above. Two differences explain the
reversal --- the powered run uses the measured hourly demand profile
rather than a daily-average magnitude, and the peak-hour topology loads
the corridor far harder --- and the direction of the correction is the
uncomfortable one: under realistic peak demand, a shielded stock SLM
makes the A501 corridor worse, not better. The fine-tuned arms
(\S\ref{sec:res-closedloop}) do not share this deficit, which is the
strongest practical argument in this dissertation for distillation.

The replication also yielded a methodological finding I had not planned
to measure. On the 26 seeds shared between runs, the \emph{catalogue}
Qwen3-0.6B artifact and my identically-prompted, identically-shielded
pipeline compile of the same weights differ by nine points of mean
relative delay ($+4.1\%$ vs $-5.2\%$ against the floor) on this
high-variance topology --- quantisation and packaging provenance alone
shift the headline sign. This is precisely why every fine-tuning
comparison in this dissertation uses the identical-pipeline compile as
its control (\S\ref{sec:met-baselines}), and it is a caution for the
field: ``the same model'' is not the same controller across serving
artifacts.

\section{Offline distillation: saturation at 0.6B}\label{sec:res-offline}
Table~\ref{tab:offline-full} gives the complete offline campaign: both
student scales, stock versus fine-tuned through the identical int4 pipeline,
plus the LOTO and format-specialist controls, evaluated on the frozen
5{,}580-state holdout as served (WebGPU, int4).

\begin{table}[t]
\caption{Holdout accuracy (\%) against frozen frontier-teacher labels,
int4-as-served, every cell backed by a committed evaluation artifact. All
fine-tuned rows and the stock Phi row emit 100\% strict JSON; stock
Qwen3-0.6B reaches 93.4\% on sota.}
\label{tab:offline-full}
\centering
\small
\begin{tabular}{lcccc}
\toprule
 & myopic & sota & prediction & coordination \\
\midrule
Qwen3-0.6B stock & 45.6 & 55.1 & 63.3 & 69.3 \\
\textbf{Qwen3-0.6B generalist ft} & \textbf{87.9} & \textbf{97.7} & \textbf{97.6} & \textbf{99.0} \\
Phi-4-mini stock & 87.5 & 59.0 & 86.0 & 87.0 \\
\textbf{Phi-4-mini generalist ft} & \textbf{88.7} & \textbf{97.1} & \textbf{97.4} & \textbf{99.3} \\
Phi-4-mini LOTO & 88.3 & 96.8 & 97.4 & 99.1 \\
Phi-4-mini sota-specialist & 87.9 & 97.8 & 89.5 & 93.3 \\
Phi-4-mini myopic-specialist & 87.0 & 87.9 & 88.8 & 91.8 \\
Phi-4-mini prediction-specialist & 87.9 & 91.3 & 97.3 & 98.4 \\
Phi-4-mini coordination-specialist & 87.4 & 89.1 & 94.7 & 99.3 \\
\bottomrule
\end{tabular}
\end{table}

Five findings follow.
\begin{enumerate}
\item \finding{Scale saturation.} The 0.6B and 3.8B generalists are
statistically indistinguishable on every format: the largest gap is
0.8 points (myopic), against a standard error of about 1.2 points on
$n{=}1{,}510$, so no two-proportion test approaches significance. Both sit
at the teacher-agreement ceiling on three of four formats (97--99\%
against 98--99\% teacher self-consistency); the myopic format is the
exception at 88\%, for the reason given below. Distillation saturates
this decision task at 0.6B; six times the parameters buy nothing
per-decision.
\item \finding{Where the gain lands depends on the base.} The 0.6B is lifted
between $+29.7$ (coordination) and $+42.3$ (sota) points across the four
formats. Phi-4-mini, already 86--88\% on three formats stock, gains most
on the delay-aware sota format ($+38.1$ points), materially on prediction
($+11.4$) and coordination ($+12.3$), and negligibly on myopic ($+1.2$).
For a capable base, fine-tuning buys decision quality, not
format discipline --- stock Phi is already 100\% format-compliant.
\item \finding{Generalists only.} No format specialist beats the generalist
on its own format by more than $0.7$ points, while off-format costs run
from $0.8$ to $9.2$ points;
and the myopic specialist loses \emph{even on-format} (87.0 vs 88.7),
consistent with mixed-format training regularising against the noisy
privileged-teacher labels that pure-myopic training overfits. This comparison was run at
3.8B only --- no Qwen format specialists were trained --- so it is a
single-scale result, unlike the generalist comparison above.
\item \finding{No topology memorisation.} The LOTO student, trained with all
Old Street states removed, matches the generalist on every format.
\item \finding{Quantised serving is exact for this task.} CPU and WebGPU
int4 artifacts agree bit-identically (97.42\% both on the sota holdout);
I scope ``lossless'' strictly as bit-identical holdout accuracy under
weight-only int4 RTN.
\end{enumerate}
The myopic ceiling ($\sim$88\%) is intrinsic: the teacher was privileged
with waiting-time state the myopic prompt hides, an Orca-style label-only
distillation cost \cite{mukherjee2023orca} disclosed as a design property,
not noise.

\section{Closed-loop outcome: delay decouples from accuracy}\label{sec:res-closedloop}
Table~\ref{tab:closedloop-full} reports the Qwen3-0.6B closed-loop sweep:
fine-tuned (ft) versus identically-compiled stock, three topologies
$\times$ 30 training-disjoint seeds, paired per-seed against fixed-time.

\begin{table}[t]
\caption{Closed-loop delay vs fixed-time and decision composition
($n{=}30$ disjoint seeds per cell; paired $t$). Negative $\Delta$ favours
the controller. \emph{auth} = share of all actuations authored by the SLM;
\emph{a-s} = share forced by the anti-starvation floor; \emph{acc} =
proposal acceptance; \emph{div} = divergence from MaxPressure among
SLM-authored decisions. The substitution path accounts for
$\le$0.5\% everywhere and is omitted.}
\label{tab:closedloop-full}
\centering
\small
\begin{tabular}{llrrrrrr}
\toprule
Topology & Arm & $\Delta$delay & $p$ & auth & a-s & acc & div \\
\midrule
Euston peak-hour & ft & $-3.9\%$ & .16 & 0.64 & 0.36 & 0.70 & 0.12 \\
 & stock & $-6.2\%$ & .12 & 0.62 & 0.38 & 0.68 & 0.47 \\
Bloomsbury grid & ft & $+0.7\%$ & .37 & 0.70 & 0.30 & 0.73 & 0.18 \\
 & stock & $+2.6\%$ & $1.5{\times}10^{-3}$ & 0.62 & 0.38 & 0.65 & 0.46 \\
Old Street & ft & $-3.1\%$ & .20 & 0.56 & 0.44 & 0.90 & 0.11 \\
 & stock & $-1.1\%$ & .35 & 0.39 & 0.61 & 0.54 & 0.60 \\
\midrule
Pooled ($n{=}90$) & ft & $-2.1\%$ & .14 & & & & \\
 & stock & $-1.6\%$ & .39 & & & & \\
\midrule
Euston & phi ft & $-3.8\%$ & .15 & 0.64 & 0.36 & 0.70 & 0.13 \\
Bloomsbury & phi ft & $+0.6\%$ & .42 & 0.70 & 0.30 & 0.73 & 0.17 \\
Old Street & phi ft & $-1.0\%$ & .56 & 0.57 & 0.43 & 0.88 & 0.12 \\
Pooled ($n{=}90$) & phi ft & $-1.4\%$ & .20 & & & & \\
\bottomrule
\end{tabular}
\end{table}

The honest headline: at $n{=}30$ per topology, neither arm significantly
beats the fixed-time floor anywhere, and a $+42$-point offline accuracy gain
produces no significant delay change --- the \finding{delay-decoupling
paradox} anticipated by the decision battery. The one significant effect of
fine-tuning is protective: stock control significantly \emph{harms} the
congested Bloomsbury grid ($+2.6\%$, $p{=}1.5{\times}10^{-3}$), the
fine-tuned controller is neutral there, and the head-to-head paired
difference favours fine-tuning by $-1.8\%$ ($p{=}.019$). I return to the
mechanism --- the shield bounds worst-case behaviour for both arms, and
external calibration puts the entire LLM-vs-MaxPressure headroom at 1--2\%
\cite{dglight2026} --- in \S\ref{sec:dis-decoupling}.
Reporting only paired $t$-tests on means, as the table above does, hides
structure that matters. The per-seed distributions are heavily skewed by
the catastrophic tail quantified in \S\ref{sec:res-forensics}, and a mean
is the wrong summary for a skewed distribution. Table~\ref{tab:robust}
therefore reports 95\% confidence intervals alongside a Wilcoxon
signed-rank test, which asks whether the controller helps on a
\emph{typical} seed rather than on average.

\begin{table}[h]
\caption{Closed-loop effects with 95\% CI on the mean relative change and a
Wilcoxon signed-rank test on the paired per-seed delays ($n{=}30$).
Bold marks $p<.05$.}
\label{tab:robust}
\centering
\small
\begin{tabular}{llrrr}
\toprule
Arm & Topology & rel.\ (95\% CI) & $p_t$ & $p_{\text{Wilcoxon}}$ \\
\midrule
Qwen ft & Euston & $-3.9$ $[-12.1, +4.4]$ & .155 & .544 \\
 & Bloomsbury & $+0.7$ $[-0.7, +2.0]$ & .371 & .360 \\
 & Old Street & $-3.1$ $[-9.0, +2.8]$ & .202 & \textbf{.011} \\
Qwen stock & Euston & $-6.2$ $[-17.6, +5.3]$ & .117 & \textbf{.018} \\
 & Bloomsbury & $+2.6$ $[+1.2, +4.1]$ & \textbf{.002} & \textbf{.003} \\
 & Old Street & $-1.1$ $[-4.4, +2.1]$ & .351 & .318 \\
Phi ft & Euston & $-3.8$ $[-11.7, +4.1]$ & .151 & .309 \\
 & Bloomsbury & $+0.6$ $[-0.7, +2.0]$ & .418 & .318 \\
 & Old Street & $-1.0$ $[-7.0, +5.1]$ & .560 & .178 \\
\bottomrule
\end{tabular}
\end{table}

Two cells change interpretation. On Old Street the fine-tuned Qwen
controller \finding{significantly beats the fixed-time floor on the
typical seed} ($p{=}.011$) while its mean effect is not significant
($p{=}.202$); the stock controller shows the same pattern on Euston
($p{=}.018$ against $p{=}.117$). The confidence intervals explain why ---
they are wide and straddle zero precisely on the topologies with the
heaviest failure tails. The honest characterisation of this controller is
therefore not ``no effect'' but \emph{typical-case improvement with a
catastrophic tail}: it helps on most seeds and fails badly on a few, and
those few are enough to erase the mean. That is a materially different
engineering problem from a controller that is simply mediocre, and it
points at tail control --- a hard per-decision timeout, or an outlier
detector that reverts to the classical policy --- rather than at better
average decisions.

I refrain from claiming these rank-test results as confirmatory findings.
Across the dissertation roughly forty paired comparisons are reported;
the two that reach $p<.05$ only under the rank test would not survive a
Holm correction applied across that family, and I did not pre-register
either. They are reported as exploratory evidence about distributional
shape, and the pre-registered comparison --- the head-to-head between
fine-tuned and stock arms --- remains the confirmatory one.

The Phi-4-mini confirmation arm (final rows of
Table~\ref{tab:closedloop-full}) extends the saturation finding
end-to-end: the 3.8B fine-tuned student's closed-loop profile is
arm-for-arm indistinguishable from the 0.6B's --- Euston $-3.8\%$ vs
$-3.9\%$, Bloomsbury $+0.6\%$ vs $+0.7\%$, the same divergence band
(0.11--0.18) and authorship profile --- and at the lowest served latency of
any arm (median p50 0.94\,s, worst cell p99 1.46\,s). The dissertation's scale-threshold
question is thereby answered in both evaluation regimes: 0.6B suffices.

\section{Demand calibration: the result the gridlock was hiding}\label{sec:res-calib}
The guard metrics of \S\ref{sec:res-guards} showed that the Bloomsbury map
was oversaturated --- under 15\% of vehicles completed and roughly 120 per
run were teleported out of deadlock. Because that map carried the only
significant closed-loop result in the study, I re-examined it rather than
merely caveating it, and the re-examination changed the dissertation's
central empirical claim.

Two diagnostics established the cause. Tripling the horizon to 3{,}600
steps raised completion only to 35\% while multiplying teleports roughly
sevenfold (121 to 830): the network was not under-simulated, it was over
capacity. Subsampling demand located the boundary sharply --- at 40\% of
the original demand the network clears (87.6\% completion, zero
teleports), while at 55\% it collapses again (57.5\%, teleports return).
I therefore added a demand-calibrated variant at 40\% (686 vehicles),
applying the same treatment Old Street had received earlier in the project
for the same reason, and re-ran all three closed-loop arms over the
identical thirty disjoint seeds, retaining the oversaturated cells.

The reader is entitled to ask whether this is a principled correction or a
search for a more favourable test bed, and the honest answer has parts on
both sides. In favour: the guard metrics that condemned the map were
committed in \S\ref{sec:met-metrics} before this analysis; the diagnosis
is independent of the outcome, since a longer horizon degraded clearance
further, which indicates demand above capacity rather than insufficient
time; the 40\% target was set to match the completion rate of the
already-calibrated Old Street map rather than chosen for its effect size;
and the same treatment had been applied to Old Street earlier in the
project for the same reason. Against: the decision to examine the map at
all followed the observation that its result was weak, which makes it a
data-dependent analysis choice, and retaining the superseded cells does
not by itself make the choice of which to headline principled.

What settles it is whether the threshold is applied uniformly. Euston
peak-hour also fails the completion guard (44--50\%), and it is the
topology on which SLM control looks \emph{worst}
(\S\ref{sec:res-factorial}), so calibrating it risks the conclusion
rather than protecting it. \todo{report the Euston calibration outcome
here whichever way it falls; until then the 18.3\% result is conditional
on a single recalibrated map rather than a general effect size.}

\begin{table}[t]
\caption{The same controllers under two demand regimes on the same
Bloomsbury network ($n{=}30$ paired seeds each). Positive $\Delta$ is
worse than the fixed-time floor.}
\label{tab:calib}
\centering
\small
\begin{tabular}{llrrrr}
\toprule
Regime & Arm & $\Delta$delay (95\% CI) & $p$ & Compl. & Auth. \\
\midrule
Oversaturated & qwen ft & $+0.7$ $[-0.7, +2.0]$ & .371 & 14.2\% & 0.70 \\
(14\% clears) & qwen stock & $+2.6$ $[+1.2, +4.1]$ & .002 & 13.5\% & 0.62 \\
\midrule
Calibrated & qwen ft & $+0.00$ $[-0.54, +0.54]$ & .987 & 87.4\% & 0.77 \\
(87\% clears) & phi ft & $+0.03$ $[-0.41, +0.46]$ & .933 & 87.5\% & 0.77 \\
 & qwen stock & $\mathbf{+22.7}$ $[+20.7, +24.6]$ & $\mathbf{6{\times}10^{-20}}$ & 76.4\% & 0.59 \\
\bottomrule
\end{tabular}
\end{table}

\finding{The gridlock was masking the effect, not creating it.} On the
calibrated network the stock controller degrades mean delay by
\risk{22.7\%} against the fixed-time floor --- an order of magnitude
larger than the $+2.6\%$ measured on the same map when it was
oversaturated --- while both fine-tuned students are statistically
indistinguishable from the floor. The head-to-head is the largest and most
certain effect in this dissertation: fine-tuning reduces delay by
\finding{$18.3\%$} relative to the identically-compiled stock model
($[-19.7, -16.9]$, $p{=}8.5{\times}10^{-19}$), and it \finding{wins on
30 of 30 seeds}. The Phi student reproduces this exactly ($-18.3\%$,
30/30, $p{=}4.5{\times}10^{-19}$).

The mechanism is straightforward once stated. A deadlocked network is
already at its worst: there is a ceiling on how much further a poor
controller can degrade it, and the teleport mechanism truncates precisely
the worst journeys, so bad control and good control are compressed toward
the same terrible mean. A free-flowing network has room for a poor
controller to do real damage --- and room for a good one to avoid it. The
stock arm's throughput corroborates this reading: it completes 76.4\% of
vehicles against 87.4\% for the fine-tuned arm, so it is stranding traffic
as well as delaying it, and its decision profile is unchanged (484 valid
proposals of 546 decisions, no fallbacks, normal latency), which rules out
a degraded-serving explanation.

\finding{This also converts the scale-saturation result from a
non-rejection into a demonstrated equivalence.} On the calibrated map the
3.8B and 0.6B students differ by $+0.03\%$ with a 95\% confidence interval
of $[-0.30, +0.36]$ --- an interval narrow enough to exclude any
practically meaningful difference, rather than merely failing to detect
one. Six times the parameters buy nothing, and now the claim is bounded
rather than assumed.

I draw a methodological conclusion beyond this project. \finding{Evaluating
signal controllers on oversaturated networks systematically understates
effect sizes.} Had I reported only the original Bloomsbury cells, this
dissertation would have concluded that fine-tuning yields a marginal
$1.8\%$ improvement of doubtful significance; the same experiment on a
network that clears yields $18.3\%$ on every seed. Since the surrounding
literature evaluates predominantly on synthetic grids under demand that is
rarely reported in clearance terms (\S\ref{sec:bg-rl}), and since the
field's own calibration figure for the whole LLM-versus-MaxPressure gap is
1--2\% \cite{dglight2026}, the possibility that published effect sizes are
compressed by saturation deserves attention. Reporting completion and
teleport counts alongside delay costs nothing and would settle it.

\section{Throughput guards: which delay numbers are comparable}\label{sec:res-guards}
Mean delay can be improved by stranding vehicles, so \S\ref{sec:met-metrics}
commits to reporting throughput, teleports and undeparted backlog alongside
it. Table~\ref{tab:guards} does so, and it changes how two of the three
topologies should be read.

\begin{table}[h]
\caption{Guard metrics, means over $n{=}30$ seeds. ``Compl.'' is completed
vehicles as a share of those loaded within the 1{,}200-step horizon;
teleports are SUMO's gridlock-escape removals.}
\label{tab:guards}
\centering
\small
\begin{tabular}{llrrrr}
\toprule
Topology & Arm & Loaded & Compl. & Teleports & Undep. \\
\midrule
Euston peak-hour & baseline & 780 & 44.3\% & 33.8 & 225 \\
 & qwen ft & 780 & 48.7\% & 28.1 & 200 \\
 & qwen stock & 780 & 50.3\% & 23.3 & 199 \\
Bloomsbury grid & baseline & 1715 & 14.6\% & 118.5 & 335 \\
 & qwen ft & 1715 & 14.2\% & 120.0 & 341 \\
 & qwen stock & 1715 & 13.5\% & 135.6 & 348 \\
Old Street & baseline & 480 & 85.5\% & 1.5 & 0 \\
 & qwen ft & 480 & 86.7\% & 0.5 & 0 \\
 & qwen stock & 480 & 84.4\% & 1.0 & 1 \\
\bottomrule
\end{tabular}
\end{table}

\finding{Only Old Street is measured in a free-flowing regime.} There 85\%
of vehicles complete and teleports are effectively absent, so mean delay is
a well-posed comparison --- and it is also the topology where the
fine-tuned controller significantly beats the floor on the typical seed
(Table~\ref{tab:robust}). The cleanest regime yields the clearest result,
which is reassuring rather than coincidental.

\risk{Bloomsbury is a gridlock-regime measurement and must be scoped as
one.} Under 15\% of vehicles complete within the horizon and over a
hundred are teleported out of deadlock, so ``mean network delay'' there is
computed over a network that never clears, with an unknown fraction of the
worst journeys removed by the simulator. I therefore treat the one
statistically significant closed-loop result in this dissertation --- the
stock controller's $+2.6\%$ harm --- as regime-scoped. It does survive the
guard test in direction: the stock arm completes fewer vehicles than the
baseline (231 vs 251) and triggers more teleports (135.6 vs 118.5), so the
delay harm is corroborated rather than contradicted by throughput, and the
fine-tuned arm sits between them on both. But the underlying comparison is
between two congested-failure modes, not between two free-flowing
policies, and a reader should weight it accordingly. Extending the horizon
or lowering demand on this map is the obvious remedy and is the first item
of unfinished work I would prioritise.

Euston sits in between at 44--50\% completion. There the SLM arms
\emph{complete more vehicles than the baseline} (380 and 392 against 345)
while showing no significant delay advantage --- a genuine
throughput-alongside-delay improvement that a delay-only report would have
missed entirely, and the reason both metrics are mandated in the protocol.

\section{Authorship: what fine-tuning actually moves}\label{sec:res-authorship}
The structural columns of Table~\ref{tab:closedloop-full} carry the
dissertation's central finding, and reading them correctly requires
keeping the three denominators of \S\ref{sec:des-shield} apart --- an
error I made in an earlier draft and correct here explicitly, because the
distinction changes the interpretation.

First, the fact that frames everything else: \finding{the shield's
substitution path almost never fires}. Across every arm and topology it
accounts for $\le$0.5\% of decisions, because both stock and fine-tuned
models emit a well-formed, in-menu proposal on essentially every tick.
Whatever the safety scaffold contributes here, it is not the correction of
malformed output. The real competition for authorship is between the SLM
and the anti-starvation floor.

Second, \finding{fine-tuning does raise the SLM's true authorship of
actuations}, and most where the stock model was weakest: on Old Street from
0.39 to 0.56, on Bloomsbury from 0.62 to 0.70, on Euston from 0.62 to 0.64.
The mechanism is visible in the adjacent column: the anti-starvation floor
is forced to intervene far less often (Old Street 0.61 $\to$ 0.44), because
the distilled model serves neglected approaches before the fairness timer
must. Proposal acceptance rises correspondingly (0.54 $\to$ 0.90). But the
honest headline number is the authored share, not the acceptance rate:
even the best fine-tuned arm leaves \emph{30--44\% of actuations authored
by a hard-coded timer}, which any audit of this deployment must attribute
to the scaffold rather than to the model.

\begin{figure}[t]
\centering
\begin{tikzpicture}
\begin{axis}[
  ybar, bar width=5.5pt, width=0.95\linewidth, height=52mm,
  ymin=0, ymax=1.0, ylabel={\sffamily\small rate},
  symbolic x coords={Euston, Bloomsbury, Old Street},
  xtick=data, enlarge x limits=0.25,
  legend style={font=\sffamily\scriptsize, at={(0.5,1.02)}, anchor=south,
    legend columns=2, draw=none},
  tick label style={font=\sffamily\small},
  ymajorgrids, major grid style={nbgrey!25},
  axis line style={nbgrey}]
\addplot[fill=nbteal!70, draw=nbteal] coordinates
  {(Euston,0.62) (Bloomsbury,0.62) (Old Street,0.39)};
\addplot[fill=nbteal!25, draw=nbteal] coordinates
  {(Euston,0.64) (Bloomsbury,0.70) (Old Street,0.56)};
\addplot[fill=nbred!70, draw=nbred] coordinates
  {(Euston,0.47) (Bloomsbury,0.46) (Old Street,0.60)};
\addplot[fill=nbred!25, draw=nbred] coordinates
  {(Euston,0.12) (Bloomsbury,0.18) (Old Street,0.11)};
\legend{authored (stock), authored (ft), divergence (stock), divergence (ft)}
\end{axis}
\end{tikzpicture}
\caption{What fine-tuning moves ($n{=}30$ seeds per bar). Teal: the SLM's
share of \emph{all} served actuations rises, most sharply on Old Street
(0.39 $\to$ 0.56); the remainder is authored by the anti-starvation floor.
Red: divergence from MaxPressure among SLM-authored decisions falls
roughly four-fold --- the distilled policy coincides with the classical
reference far more often.}
\label{fig:authorship}
\end{figure}

Third, and cutting against the simple reading, \finding{the distilled
policy is markedly more MaxPressure-like}: divergence among SLM-authored
decisions falls from 0.46--0.60 to 0.11--0.18, so the fine-tuned model
chooses what MaxPressure would have chosen on roughly nine decisions in
ten. This admits a deflationary rival explanation --- that distillation
largely taught the student to imitate the classical reference --- and the
evidence neither establishes nor excludes it. Three observations bound it.
The teacher was a frontier model reasoning over accumulated waiting, not
MaxPressure, and the students agree with the teacher's held-out labels at
97--99\% (\S\ref{sec:res-offline}), so the convergence is at most
indirect. The fine-tuned arms do not reproduce MaxPressure's closed-loop
behaviour --- MaxPressure is significantly worse than fixed-time on these
topologies (\S\ref{sec:res-maxpressure}) while the fine-tuned arms are
not. And the remaining one-in-ten divergences are where the arms differ
from both baselines. Distinguishing imitation from independent
rediscovery of a similar policy would require a decision-level
counterfactual study I did not run, and I flag it as the most substantive
open question about this result.

What survives all three qualifications is the audit-relevant claim, stated
at its true magnitude. An accident audit over the stock Old Street
deployment could attribute only 39\% of actuations to the learned
component; over the fine-tuned deployment it attributes 56\%, with the
balance falling to a deterministic timer whose behaviour is trivially
explicable. Fine-tuning buys a controller whose served decisions are
substantially its own and which needs the fairness scaffold less --- a
smaller claim than \emph{the model is the controller}, and one the data
supports.

\section{Failure forensics: the audit layer explains the tail}\label{sec:res-forensics}
The audit architecture's value shows most clearly when control fails. Across
all 270 closed-loop cells, 14 (5.2\%) are catastrophic (delay
$\ge{+}20\%$ vs baseline) --- and every one of their signed decision chains
verifies, so each failure has a provable record to interrogate. Three
findings emerge from that interrogation.

First, \finding{failures are a controller property, not seed hardness}: the
catastrophic seeds are baseline-easy or median demand realisations (baseline
delay ranks 1--20 of $\sim$35, never the hard tail), and no seed breaks all
three arms.

Second, the two failure regimes have different, attributable mechanisms. The
stock model's five failures are its own served decisions: in the worst cell
(Euston s13, $+110.3\%$), the verified chain shows 214 of 335 decisions
SLM-authored with 119 divergent proposals served (override 0.556), zero
shield fallbacks, and anti-starvation at its normal rate --- the scaffold
did not cause the collapse, the model's policy did, and the chain proves it.
The same-seed fine-tuned arms sit at $+1.1\%$ and $+5.5\%$.

Third, \finding{the two fine-tuned students converge at the trajectory
level on a substantial minority of cells}. Comparing the 0.6B and 3.8B
students seed-for-seed across the 90 shared cells: on nine they produce
\emph{bit-identical} mean delay, on a further 25 they agree within 0.5
points, and on the remaining 56 they differ. The identical cells are not
duplicated records --- wall-clock times differ by up to a factor of two
and median latencies by 2.1$\times$ --- but genuinely separate runs that
served the same phase sequence; in one such cell both arms took 307
decisions and their SLM/floor attribution differs by exactly one tick,
where the floor and the model happened to select the same phase. The
catastrophic cells fall inside this converged subset, which is why the two
students fail together.

Two independently trained students, different in architecture and
$6\times$ apart in scale, therefore reproduce each other's control
trajectory on roughly a third of cells and diverge on the rest. I read
this as the behavioural counterpart of the offline saturation result ---
the distillation target appears to constrain not only what the students
score but what they \emph{do} --- while stating the limit of the evidence:
partial trajectory convergence is equally consistent with both students
defaulting to similar behaviour on the easier cells, and I have not
separated those explanations.

One honest gap: this campaign persisted aggregate decision statistics per
cell, not per-step audit records, so \emph{when} within a run the
divergence concentrated is unrecoverable here; the per-step chain is
retained in the dedicated accident-suite runs (\S\ref{sec:res-audit}).
Full tables in \texttt{results/failure\_forensics.md}.

\section{Audit and crash reconstruction}\label{sec:res-audit}
The audit layer was evaluated on seven pre-defined collision scenarios
spanning the legally distinct cases (red-light running, EV exemption,
maintenance vehicle, amber gambling, conflicting green, dark signal, spoofed
EV). All seven prove end-to-end: hash-chain verification passes, signature
verification passes, and an imposter re-signing attack is caught in every
scenario. The governing UK rule is \emph{inferred deterministically from the
verified facts} --- e.g.\ the spoofed-EV scenario correctly denies the
exemption and anchors fault on the driver, while the conflicting-green
scenario is the sole case attributing (weak, misfeasance-only) authority
fault --- and the system emits a cited evidence pack for a human
adjudicator, never a binding verdict, by design
\cite{dahl2024legal}.

\subsection{A worked evidence pack}\label{sec:res-audit-worked}
Figure~\ref{fig:evpack} traces the hardest of the seven scenarios end to
end, because it is the one where the \emph{system itself} may be at
fault: a collision at a junction that emitted a conflicting green.

Three properties are worth drawing out. First, the pack is built only
from facts the chain verifies --- \texttt{signal\_state}, the
\texttt{conflicting\_green} flag, the crossing class --- and each cited
rule carries the specific fact that triggered it, so an investigator can
audit the inference, not merely read a conclusion. Second, the legal
reasoning is faithful to a doctrine that is unfavourable to the deployer:
under \emph{Gorringe v Calderdale} a highway authority's mere omission
attracts no duty, so authority fault defaults to zero, and it becomes
non-zero here only because \emph{Bird v Pearce} treats a positively-wrong
indication as misfeasance --- fault weight \texttt{low}, flagged
\texttt{common\_law} rather than \texttt{MUST}, with legal causation
explicitly \texttt{NOT\_ASSESSED}. The system does not overstate a case
against a driver when its own record shows the signal misled them, and it
does not overstate the case against itself either. Third, the tamper
probe passes: an imposter re-signing the incident record is rejected on
signature verification even though the chain hashes are recomputed
consistently --- the property that makes the pack evidence rather than
assertion.

\begin{figure}[t]
\small
\begin{lstlisting}
VERIFIED RECORD (seq 0, junction J, hash a0bb851c..., t=100.0)
  chain verified: True   signatures verified: True
  imposter re-sign attempt: CAUGHT
VERIFIED FACTS
  signal_state = green      conflicting_green = True
  crossing_class = civilian second_party_on_green = False
  key_present = False       corroborated = False
GOVERNING RULES (inferred from the facts above)
  LR-authority-misfeasance [Bird v Pearce]  fault: low  (common law)
    "the audit shows the system emitted a conflicting green
     (a positively-wrong indication) -- misfeasance, not omission"
  LR-authority-nonfeasance [Gorringe v Calderdale] fault: low (common law)
    "authority fault otherwise defaults to zero; it turns non-zero only
     on this misfeasance/conflicting-green branch (low confidence;
     legal_causation NOT_ASSESSED)"
OUTPUT: cited evidence pack for a human adjudicator -- not an adjudication
\end{lstlisting}
\caption{The conflicting-green evidence pack, verbatim from
\texttt{results/experiment\_crash\_law.json}: the only one of the seven
scenarios in which authority fault is non-zero, and the case where the
system's own record is used against its operator.}
\label{fig:evpack}
\end{figure}

\section{Trust-coefficient battery}\label{sec:res-trust}
The trust mechanism passes its full property battery (9/9): a persistently
honest junction climbs to 0.929 trust; a persistent liar collapses to 0.011
and loses the power to corroborate; the asymmetry is deliberate (ten
verified truths to climb from 0.5 to 0.9, one lie from high trust falls
0.65); and a truthful-but-late report is vindicated retroactively by local
sensing, so honesty is never permanently penalised for latency. Trust is
discount-only by construction --- it can strip a liar's corroboration power
but can never add a preemption on zero evidence, and local sensing is never
itself doubted.

In closed-loop traffic (honest / blatant / stealth misreporting injected
into SUMO), I ran the battery against \emph{both} stock and fine-tuned
controllers at both scales --- the fine-tuned deployment is the one that
matters, precisely because of \S\ref{sec:res-authorship}: a controller that
actually acts on its coordination inputs is the one poisoned inputs can
move. Three findings (three seeds per cell, descriptive).

First, \finding{detection is controller-independent by design and in
measurement}: the ledger scores claims against local sensing, and under
every controller a blatant liar is caught and collapsed (Bloomsbury: 62--70
catches, final trust 0.21--0.27; grid4$\times$4: 146--151 catches, trust
$\le$0.03, corroboration power stripped), while a calibrated stealth
attacker evades essentially everywhere ($\le$2 catches, trust $\ge$0.92) ---
the stealth evasion stands as an open limitation. On the corridor and Old
Street the liar junction is never consulted, so the mechanism is inert
there (zero lies injectable) --- reported, not hidden.

Second, \finding{the delay cost of blatant lying survives detection}:
on the synthetic grids blatant misreporting adds roughly $+3$ to $+11$
points of delay relative to the same controller's honest run even though
the liar's trust collapses --- trust bounds the liar's future influence
but does not undo congestion already caused. On Bloomsbury, the real
grid, local-sensing override holds the damage to within $\sim$3 points
(fine-tuned qwen: $-1.3\%$ under blatant attack vs $-1.1\%$ honest).

Third, \finding{stealth's harm is positional, not congestive}: the
stealth attacker's small, intermittent inflations cost approximately
nothing in delay ($-2$ to $+3$ points vs honest) --- what stealth buys is
retained corroboration power at high trust, a position it could later
abuse. The threat model that matters for a fine-tuned deployment is
therefore asymmetric: blatant attacks are caught but briefly expensive on
congested grids; stealth attacks are cheap, patient, and undetected. Both
conclusions argue that the trust ledger is a necessary companion to
fine-tuning rather than an optional extra, and that detection calibrated
to distributional implausibility (\S\ref{sec:dis-future}) is the right
next defence.

\section{Emergency preemption: what the priority path is worth}\label{sec:res-ev}
The audit and authorisation machinery exists to serve a concrete function,
and this section measures it. Three modes are compared on the corridor over
thirty seeds, with a scripted ambulance and a scripted spoofing attack:
\emph{nopreempt}, which grants emergency vehicles no signal priority at all
and is what a conventional fixed-time installation does; \emph{naive},
which preempts on any claim of emergency status; and \emph{defended}, which
preempts only on a claim carrying a valid authorising key and corroboration
(\S\ref{sec:des-trust}).

\begin{table}[h]
\caption{Emergency-vehicle priority, $n{=}30$ seeds, 95\% CIs. Positive
values are seconds saved.}
\label{tab:ev}
\centering
\small
\begin{tabular}{lrr}
\toprule
Quantity & Mean & 95\% CI \\
\midrule
Ambulance time saved by priority (nopreempt $-$ defended) & 31.1\,s & $[25.4, 36.5]$ \\
Cost to the ambulance of the authorisation gate (defended $-$ naive) & 8.0\,s & $[3.7, 12.5]$ \\
Cross-traffic harm avoided under attack, mean & 4.6\,s & $[3.3, 6.0]$ \\
Cross-traffic harm avoided under attack, worst case & 16.1\,s & $[11.5, 20.4]$ \\
\bottomrule
\end{tabular}
\end{table}

Three functional readings follow, and all four intervals exclude zero.
\finding{Signal priority is worth about half a minute to an emergency
vehicle} on this corridor: 31.1\,s against a controller that offers no
priority. \finding{Security is not free, and its price is measurable}: the
authorisation gate costs the ambulance 8.0\,s relative to a naive
controller that preempts on any claim, because a claim must be
authenticated and corroborated before it is honoured.
\finding{That price buys something quantified}: when a spoofed emergency
claim is injected, the gate spares cross-traffic 4.6\,s of mean delay and
16.1\,s in the worst case, harm the naive controller absorbs in full.
Whether 8\,s of ambulance time is worth 4.6\,s of mean cross-traffic
protection is a policy judgement rather than an engineering one, and the
contribution here is that the exchange rate is measured rather than
asserted.

Two scope limits. The scenario is scripted, with one ambulance and one
attack window per run, so these are per-event magnitudes and not an
expected benefit under a naturalistic emergency-call arrival rate. And the
comparison is against a no-priority baseline, which characterises a
conventional installation without emergency-vehicle detection rather than
any specific deployed system.

\section{The fair re-test of the SLM authorisation negative}\label{sec:res-lee}
An early measured negative --- the raw SLM losing to a deterministic
template on its legal-note job --- carried a supervisor-posed objection
(Stott): the model was guessing because the policy was never in front of
it, so the negative cannot stand until the SLM is re-tested with explicit
in-prompt rules and a reason-then-classify structure. I ran that re-test
on the balanced 26-claim authorisation dataset of \S\ref{sec:res-audit}'s
suite, four models $\times$ two arms (policy-scaffolded vs bare), scored
against the deterministic knowledge-base gate.

The objection is partially vindicated: for a capable stock model the
scaffold works. Stock Phi-4-mini rises from 42.3\% to 69.2\% class
accuracy, its false-negative rate falls to zero, and it stops denying any
genuine emergency. But the negative itself is robust: even the best cell
--- stock Phi, scaffolded --- still makes \risk{four dangerous false
grants} (granting preemption to claims that must be denied, FPR 0.33)
where the deterministic gate makes none in 26/26. Policy grounding
rescues recall, not safety.

The re-test also produced this dissertation's clearest evidence of
specialisation cost: the \emph{fine-tuned} students are far worse here
than their own stock bases --- ft-Phi falls to 19.2\% scaffolded with 15
of 26 responses unparseable, and ft-Qwen cannot emit the required
\texttt{\{"class", "grant"\}} schema at all bare (26/26 parse failures),
because distillation captured its output format onto
\texttt{\{"phase": N\}}. The retention probe's mild MMLU decline
(\S\ref{sec:dis-threats}) understates the cost on near-neighbour
structured tasks. Both conclusions reinforce the architecture: the
deterministic gate owns authorisation, the fine-tuned model owns phase
selection, and neither should borrow the other's job.

\section{Serving reliability at the edge}\label{sec:res-reliability}
Sustained on-device serving produced four documented reliability findings
across the campaign: (i) the catalogue Qwen3-0.6B GPU artifact emitting
corrupted tokens after long service uptime; (ii, iii) two unprompted Foundry
Local service self-restarts mid-evaluation, killing runs at 5{,}500/5{,}580
and 4{,}300/5{,}580 rows; (iv) registration of builder-compiled models
requiring undocumented synthesis of the inference manifest. Each was
mitigated in harness --- endpoint rediscovery, model reload with
retry-of-failed-rows and statistics intact, fail-loud parse accounting that
aborts rather than scores error text --- and the mitigations are themselves
contributions: they are what turned a research prototype into a harness that
survived a six-week unattended campaign.

Decision latency, measured per cell across the closed-loop sweeps, is
summarised in Table~\ref{tab:latency}. The 3.8B student is the fastest arm
served (median p50 0.94\,s) and the 0.6B students the slowest
(1.74--1.75\,s) --- a counter-intuitive ordering that is a property of the
compiled artifacts and their execution paths rather than of parameter
count, and one I report because it cuts against the assumption that the
smaller model is always the cheaper one to serve. Two individual calls in
the 0.6B arms exceeded the 10\,s decision interval (worst 13.8\,s), so the
real-time margin is comfortable at the median and not guaranteed in the
tail: on those ticks the junction acted on stale state. This is the same
criterion that excludes phi-4-mini-reasoning
(\S\ref{sec:met-baselines}); the difference is frequency --- isolated
outliers here against a p99 above the interval there --- and a deployment
would need a hard timeout with deterministic fallback, which the
architecture already supports.

\begin{table}[h]
\caption{Decision latency across closed-loop cells (per-cell summaries,
GPU-served, sota arms). The decision interval is 10\,s.}
\label{tab:latency}
\centering
\small
\begin{tabular}{lrrrr}
\toprule
Arm & p50 (median) & p50 range & p99 (worst cell) & worst call \\
\midrule
Qwen3-0.6B ft & 1.75\,s & 1.39--2.57\,s & 7.58\,s & 13.79\,s \\
Qwen3-0.6B stock & 1.74\,s & 1.46--2.51\,s & 2.97\,s & 9.88\,s \\
Phi-4-mini ft & 0.94\,s & 0.92--1.20\,s & 1.46\,s & 1.49\,s \\
\bottomrule
\end{tabular}
\end{table}
